% Options for packages loaded elsewhere
% Options for packages loaded elsewhere
\PassOptionsToPackage{unicode}{hyperref}
\PassOptionsToPackage{hyphens}{url}
\PassOptionsToPackage{dvipsnames,svgnames,x11names}{xcolor}
%
\documentclass[
  spanish,
  letterpaper,
  DIV=11,
  numbers=noendperiod]{scrreprt}
\usepackage{xcolor}
\usepackage{amsmath,amssymb}
\setcounter{secnumdepth}{5}
\usepackage{iftex}
\ifPDFTeX
  \usepackage[T1]{fontenc}
  \usepackage[utf8]{inputenc}
  \usepackage{textcomp} % provide euro and other symbols
\else % if luatex or xetex
  \usepackage{unicode-math} % this also loads fontspec
  \defaultfontfeatures{Scale=MatchLowercase}
  \defaultfontfeatures[\rmfamily]{Ligatures=TeX,Scale=1}
\fi
\usepackage{lmodern}
\ifPDFTeX\else
  % xetex/luatex font selection
\fi
% Use upquote if available, for straight quotes in verbatim environments
\IfFileExists{upquote.sty}{\usepackage{upquote}}{}
\IfFileExists{microtype.sty}{% use microtype if available
  \usepackage[]{microtype}
  \UseMicrotypeSet[protrusion]{basicmath} % disable protrusion for tt fonts
}{}
\makeatletter
\@ifundefined{KOMAClassName}{% if non-KOMA class
  \IfFileExists{parskip.sty}{%
    \usepackage{parskip}
  }{% else
    \setlength{\parindent}{0pt}
    \setlength{\parskip}{6pt plus 2pt minus 1pt}}
}{% if KOMA class
  \KOMAoptions{parskip=half}}
\makeatother
% Make \paragraph and \subparagraph free-standing
\makeatletter
\ifx\paragraph\undefined\else
  \let\oldparagraph\paragraph
  \renewcommand{\paragraph}{
    \@ifstar
      \xxxParagraphStar
      \xxxParagraphNoStar
  }
  \newcommand{\xxxParagraphStar}[1]{\oldparagraph*{#1}\mbox{}}
  \newcommand{\xxxParagraphNoStar}[1]{\oldparagraph{#1}\mbox{}}
\fi
\ifx\subparagraph\undefined\else
  \let\oldsubparagraph\subparagraph
  \renewcommand{\subparagraph}{
    \@ifstar
      \xxxSubParagraphStar
      \xxxSubParagraphNoStar
  }
  \newcommand{\xxxSubParagraphStar}[1]{\oldsubparagraph*{#1}\mbox{}}
  \newcommand{\xxxSubParagraphNoStar}[1]{\oldsubparagraph{#1}\mbox{}}
\fi
\makeatother


\usepackage{longtable,booktabs,array}
\newcounter{none} % for unnumbered tables
\usepackage{calc} % for calculating minipage widths
% Correct order of tables after \paragraph or \subparagraph
\usepackage{etoolbox}
\makeatletter
\patchcmd\longtable{\par}{\if@noskipsec\mbox{}\fi\par}{}{}
\makeatother
% Allow footnotes in longtable head/foot
\IfFileExists{footnotehyper.sty}{\usepackage{footnotehyper}}{\usepackage{footnote}}
\makesavenoteenv{longtable}
\usepackage{graphicx}
\makeatletter
\newsavebox\pandoc@box
\newcommand*\pandocbounded[1]{% scales image to fit in text height/width
  \sbox\pandoc@box{#1}%
  \Gscale@div\@tempa{\textheight}{\dimexpr\ht\pandoc@box+\dp\pandoc@box\relax}%
  \Gscale@div\@tempb{\linewidth}{\wd\pandoc@box}%
  \ifdim\@tempb\p@<\@tempa\p@\let\@tempa\@tempb\fi% select the smaller of both
  \ifdim\@tempa\p@<\p@\scalebox{\@tempa}{\usebox\pandoc@box}%
  \else\usebox{\pandoc@box}%
  \fi%
}
% Set default figure placement to htbp
\def\fps@figure{htbp}
\makeatother



\ifLuaTeX
\usepackage[bidi=basic,shorthands=off]{babel}
\else
\usepackage[bidi=default,shorthands=off]{babel}
\fi
\ifLuaTeX
  \usepackage{selnolig} % disable illegal ligatures
\fi


\setlength{\emergencystretch}{3em} % prevent overfull lines

\providecommand{\tightlist}{%
  \setlength{\itemsep}{0pt}\setlength{\parskip}{0pt}}



 


\KOMAoption{captions}{tableheading}
\makeatletter
\@ifpackageloaded{tcolorbox}{}{\usepackage[skins,breakable]{tcolorbox}}
\@ifpackageloaded{fontawesome5}{}{\usepackage{fontawesome5}}
\definecolor{quarto-callout-color}{HTML}{909090}
\definecolor{quarto-callout-note-color}{HTML}{0758E5}
\definecolor{quarto-callout-important-color}{HTML}{CC1914}
\definecolor{quarto-callout-warning-color}{HTML}{EB9113}
\definecolor{quarto-callout-tip-color}{HTML}{00A047}
\definecolor{quarto-callout-caution-color}{HTML}{FC5300}
\definecolor{quarto-callout-color-frame}{HTML}{acacac}
\definecolor{quarto-callout-note-color-frame}{HTML}{4582ec}
\definecolor{quarto-callout-important-color-frame}{HTML}{d9534f}
\definecolor{quarto-callout-warning-color-frame}{HTML}{f0ad4e}
\definecolor{quarto-callout-tip-color-frame}{HTML}{02b875}
\definecolor{quarto-callout-caution-color-frame}{HTML}{fd7e14}
\makeatother
\makeatletter
\@ifpackageloaded{bookmark}{}{\usepackage{bookmark}}
\makeatother
\makeatletter
\@ifpackageloaded{caption}{}{\usepackage{caption}}
\AtBeginDocument{%
\ifdefined\contentsname
  \renewcommand*\contentsname{Tabla de contenidos}
\else
  \newcommand\contentsname{Tabla de contenidos}
\fi
\ifdefined\listfigurename
  \renewcommand*\listfigurename{Listado de Figuras}
\else
  \newcommand\listfigurename{Listado de Figuras}
\fi
\ifdefined\listtablename
  \renewcommand*\listtablename{Listado de Tablas}
\else
  \newcommand\listtablename{Listado de Tablas}
\fi
\ifdefined\figurename
  \renewcommand*\figurename{Figura}
\else
  \newcommand\figurename{Figura}
\fi
\ifdefined\tablename
  \renewcommand*\tablename{Tabla}
\else
  \newcommand\tablename{Tabla}
\fi
}
\@ifpackageloaded{float}{}{\usepackage{float}}
\floatstyle{ruled}
\@ifundefined{c@chapter}{\newfloat{codelisting}{h}{lop}}{\newfloat{codelisting}{h}{lop}[chapter]}
\floatname{codelisting}{Listado}
\newcommand*\listoflistings{\listof{codelisting}{Listado de Listados}}
\makeatother
\makeatletter
\makeatother
\makeatletter
\@ifpackageloaded{caption}{}{\usepackage{caption}}
\@ifpackageloaded{subcaption}{}{\usepackage{subcaption}}
\makeatother
\usepackage{bookmark}
\IfFileExists{xurl.sty}{\usepackage{xurl}}{} % add URL line breaks if available
\urlstyle{same}
\hypersetup{
  pdftitle={Módulo de Instalaciones Civiles e Industriales},
  pdfauthor={Ing. Jaime Rojas C., Mgtr.},
  pdflang={es},
  colorlinks=true,
  linkcolor={blue},
  filecolor={Maroon},
  citecolor={Blue},
  urlcolor={Blue},
  pdfcreator={LaTeX via pandoc}}


\title{Módulo de Instalaciones Civiles e Industriales}
\usepackage{etoolbox}
\makeatletter
\providecommand{\subtitle}[1]{% add subtitle to \maketitle
  \apptocmd{\@title}{\par {\large #1 \par}}{}{}
}
\makeatother
\subtitle{Normalización, diagramas, instrumentación y automatización}
\author{Ing. Jaime Rojas C., Mgtr.}
\date{2026-05-28}
\begin{document}
\maketitle

\renewcommand*\contentsname{Tabla de contenidos}
{
\setcounter{tocdepth}{2}
\tableofcontents
}

\bookmarksetup{startatroot}

\chapter{Inicio}\label{inicio}

Bienvenido al \textbf{Módulo de Instalaciones Civiles e Industriales}.

Este libro digital ha sido organizado para la asignatura
\textbf{Instalaciones Civiles e Industriales}, correspondiente al cuarto
ciclo de la Carrera de Electricidad de la Universidad Católica de
Cuenca. Su finalidad es integrar fundamentos técnicos, criterios
normativos, interpretación de planos, mediciones eléctricas, seguridad,
automatización básica y prácticas experimentales.

\begin{tcolorbox}[enhanced jigsaw, arc=.35mm, bottomrule=.15mm, bottomtitle=1mm, breakable, colback=white, colbacktitle=quarto-callout-note-color!10!white, colframe=quarto-callout-note-color-frame, coltitle=black, left=2mm, leftrule=.75mm, opacityback=0, opacitybacktitle=0.6, rightrule=.15mm, title=\textcolor{quarto-callout-note-color}{\faInfo}\hspace{0.5em}{Enfoque del módulo}, titlerule=0mm, toprule=.15mm, toptitle=1mm]

El módulo sigue la estructura oficial de tres bloques:
\textbf{Normalización Eléctrica, Seguridad y Simbología};
\textbf{Diagramas Eléctricos e Instrumentación}; y
\textbf{Automatización}.

\end{tcolorbox}

\section{Resultado de aprendizaje de la
asignatura}\label{resultado-de-aprendizaje-de-la-asignatura}

Aplicar los principios de los diseños de circuitos eléctricos, su
simbología e interpretación práctica.

\section{Objetivo general del
módulo}\label{objetivo-general-del-muxf3dulo}

Desarrollar competencias para interpretar, diseñar y analizar
instalaciones eléctricas civiles e industriales, aplicando normalización
eléctrica, simbología técnica, diagramas eléctricos, mediciones
eléctricas, análisis de errores, seguridad en laboratorio y fundamentos
de automatización mediante PLC.

\section{Organización del
aprendizaje}\label{organizaciuxf3n-del-aprendizaje}

{\def\LTcaptype{none} % do not increment counter
\begin{longtable}[]{@{}lr@{}}
\toprule\noalign{}
Componente & Horas \\
\midrule\noalign{}
\endhead
\bottomrule\noalign{}
\endlastfoot
Aprendizaje en contacto con el docente & 32 \\
Aprendizaje autónomo & 32 \\
Aprendizaje práctico-experimental & 32 \\
\textbf{Total} & \textbf{96} \\
\end{longtable}
}

\section{Cómo usar este libro}\label{cuxf3mo-usar-este-libro}

Cada bloque contiene:

\begin{enumerate}
\def\labelenumi{\arabic{enumi}.}
\tightlist
\item
  Resultado de aprendizaje.
\item
  Objetivos específicos.
\item
  Desarrollo conceptual.
\item
  Aplicaciones técnicas.
\item
  Actividades de aprendizaje.
\item
  Prácticas vinculadas.
\item
  Evaluación sugerida.
\end{enumerate}

\bookmarksetup{startatroot}

\chapter{Plan de la asignatura}\label{plan-de-la-asignatura}

\section{Datos informativos}\label{datos-informativos}

{\def\LTcaptype{none} % do not increment counter
\begin{longtable}[]{@{}ll@{}}
\toprule\noalign{}
Campo & Información \\
\midrule\noalign{}
\endhead
\bottomrule\noalign{}
\endlastfoot
Asignatura & Instalaciones Civiles e Industriales \\
Carrera & Electricidad \\
Ciclo & Cuarto ciclo \\
Código & IC-IE-04 \\
Créditos & 2 \\
Modalidad & Presencial \\
Unidad de organización curricular & Básica \\
Disciplina & Ciencias Tecnológicas \\
Subdisciplina & Generación, Transmisión y Distribución \\
\end{longtable}
}

\section{Caracterización de la
asignatura}\label{caracterizaciuxf3n-de-la-asignatura}

La asignatura proporciona fundamentos teóricos y prácticos para la
interpretación, diseño y análisis de instalaciones eléctricas en
edificaciones civiles y entornos industriales. Integra normalización
eléctrica, simbología técnica, diagramas eléctricos, seguridad,
mediciones eléctricas, análisis de errores y fundamentos de
automatización industrial.

\section{Bloques formativos}\label{bloques-formativos}

{\def\LTcaptype{none} % do not increment counter
\begin{longtable}[]{@{}
  >{\raggedright\arraybackslash}p{(\linewidth - 6\tabcolsep) * \real{0.2500}}
  >{\raggedright\arraybackslash}p{(\linewidth - 6\tabcolsep) * \real{0.2500}}
  >{\raggedright\arraybackslash}p{(\linewidth - 6\tabcolsep) * \real{0.2500}}
  >{\raggedright\arraybackslash}p{(\linewidth - 6\tabcolsep) * \real{0.2500}}@{}}
\toprule\noalign{}
\begin{minipage}[b]{\linewidth}\raggedright
Bloque
\end{minipage} & \begin{minipage}[b]{\linewidth}\raggedright
Título
\end{minipage} & \begin{minipage}[b]{\linewidth}\raggedright
Fechas referenciales
\end{minipage} & \begin{minipage}[b]{\linewidth}\raggedright
Práctica asociada
\end{minipage} \\
\midrule\noalign{}
\endhead
\bottomrule\noalign{}
\endlastfoot
1 & Normalización Eléctrica, Seguridad y Simbología & 20-03-2026 al
24-04-2026 & Diseño de plano eléctrico residencial \\
2 & Diagramas Eléctricos e Instrumentación & 01-05-2026 al 05-06-2026 &
Diseño de plano eléctrico industrial \\
3 & Automatización & 12-06-2026 al 17-07-2026 & Fundamentos de
programación con PLCs \\
\end{longtable}
}

\section{Bibliografía base}\label{bibliografuxeda-base}

\begin{itemize}
\tightlist
\item
  Davie, A. G. (1974). \emph{Introducción a la automatización
  industrial}. Universitaria.
\item
  Enríquez Harper, G. (2017). \emph{Protección de instalaciones
  eléctricas industriales y comerciales}. Limusa.
\item
  Enríquez Harper, G. (2012). \emph{Guía para el diseño de instalaciones
  eléctricas residenciales, industriales y comerciales}. Limusa.
\end{itemize}

\section{Bibliografía
complementaria}\label{bibliografuxeda-complementaria}

\begin{itemize}
\tightlist
\item
  García Trasancos, J. (2011). \emph{Instalaciones eléctricas en media y
  baja tensión}. Paraninfo.
\item
  MIDUVI. (2018). \emph{NEC-SB-IE: Instalaciones eléctricas}. Norma
  Ecuatoriana de la Construcción.
\item
  Festo. (2022). \emph{Automatización industrial: sistemas de
  aprendizaje para la formación técnica}.
\end{itemize}

\bookmarksetup{startatroot}

\chapter{Bloque 1: Normalización Eléctrica, Seguridad y
Simbología}\label{bloque-1-normalizaciuxf3n-eluxe9ctrica-seguridad-y-simbologuxeda}

\section{Presentación del bloque}\label{presentaciuxf3n-del-bloque}

El Bloque 1 introduce al estudiante en los fundamentos técnicos,
normativos y gráficos necesarios para interpretar instalaciones
eléctricas civiles e industriales. En esta etapa se estudia la
importancia de la normalización eléctrica, el uso de simbología técnica,
la identificación de componentes en planos eléctricos y la aplicación de
criterios de seguridad en laboratorio y obra.

Una instalación eléctrica no debe entenderse únicamente como la conexión
de cables, interruptores, luminarias y tomacorrientes. En términos
técnicos, es un sistema integrado que permite transportar, distribuir,
controlar y utilizar la energía eléctrica de manera segura, eficiente y
conforme a normas.

En el ámbito civil, las instalaciones se orientan principalmente a
viviendas, edificios, locales comerciales, oficinas, instituciones
educativas y edificaciones de servicio. En el ámbito industrial, las
instalaciones incorporan motores, tableros de potencia, tableros de
control, sensores, actuadores, canalizaciones industriales, sistemas de
protección y tecnologías de automatización.

\begin{tcolorbox}[enhanced jigsaw, arc=.35mm, bottomrule=.15mm, bottomtitle=1mm, breakable, colback=white, colbacktitle=quarto-callout-note-color!10!white, colframe=quarto-callout-note-color-frame, coltitle=black, left=2mm, leftrule=.75mm, opacityback=0, opacitybacktitle=0.6, rightrule=.15mm, title=\textcolor{quarto-callout-note-color}{\faInfo}\hspace{0.5em}{Nota}, titlerule=0mm, toprule=.15mm, toptitle=1mm]

Este bloque constituye la base para los siguientes temas de la
asignatura: diagramas eléctricos, instrumentación, mediciones eléctricas
y automatización industrial.

\end{tcolorbox}

\section{Resultado de aprendizaje del
bloque}\label{resultado-de-aprendizaje-del-bloque}

Al finalizar este bloque, el estudiante interpreta planos eléctricos,
reconoce simbología normalizada y componentes de instalaciones
eléctricas civiles e industriales conforme a normas técnicas vigentes.
Además, aplica normas de seguridad eléctrica y protocolos de
laboratorio, identificando riesgos eléctricos y proponiendo medidas de
prevención en actividades prácticas y entornos reales.

\section{Objetivo específico del
bloque}\label{objetivo-especuxedfico-del-bloque}

Comprender y aplicar los fundamentos de las instalaciones eléctricas
civiles e industriales, incorporando la normalización eléctrica, las
normas técnicas y la simbología eléctrica, para la correcta
interpretación de planos eléctricos y la ejecución de actividades bajo
criterios de seguridad en laboratorios y entornos eléctricos.

\begin{center}\rule{0.5\linewidth}{0.5pt}\end{center}

\section{1.1 Introducción a las instalaciones eléctricas civiles e
industriales}\label{introducciuxf3n-a-las-instalaciones-eluxe9ctricas-civiles-e-industriales}

\subsection{Concepto general de instalación
eléctrica}\label{concepto-general-de-instalaciuxf3n-eluxe9ctrica}

Una instalación eléctrica es el conjunto de elementos destinados a
conducir, distribuir, proteger, controlar y utilizar la energía
eléctrica dentro de una edificación o sistema productivo.

Está conformada por:

\begin{itemize}
\tightlist
\item
  acometida;
\item
  medidor;
\item
  tablero general;
\item
  tableros secundarios;
\item
  protecciones;
\item
  conductores;
\item
  canalizaciones;
\item
  cajas de paso;
\item
  interruptores;
\item
  tomacorrientes;
\item
  luminarias;
\item
  cargas especiales;
\item
  sistema de puesta a tierra;
\item
  planos y diagramas técnicos.
\end{itemize}

En una edificación, la instalación eléctrica debe permitir que la
energía llegue desde el punto de suministro hasta cada carga de manera
segura y funcional.

\subsection{Figura 1.1. Esquema general de una instalación
eléctrica}\label{figura-1.1.-esquema-general-de-una-instalaciuxf3n-eluxe9ctrica}

\includegraphics[width=15.49in,height=4.23in]{bloque_01_normalizacion_seguridad_simbologia_files/figure-latex/mermaid-figure-1.png}

La figura muestra la secuencia general de una instalación eléctrica de
baja tensión. La energía ingresa desde la red de distribución, pasa por
la acometida, se registra en el medidor, se protege mediante un
dispositivo principal y se distribuye hacia circuitos derivados.

\subsection{Instalaciones eléctricas
civiles}\label{instalaciones-eluxe9ctricas-civiles}

Las instalaciones civiles se encuentran en edificaciones destinadas a
vivienda, comercio, educación, salud, oficinas, servicios públicos y
espacios similares.

Sus cargas más comunes son:

{\def\LTcaptype{none} % do not increment counter
\begin{longtable}[]{@{}
  >{\raggedright\arraybackslash}p{(\linewidth - 2\tabcolsep) * \real{0.5000}}
  >{\raggedright\arraybackslash}p{(\linewidth - 2\tabcolsep) * \real{0.5000}}@{}}
\toprule\noalign{}
\begin{minipage}[b]{\linewidth}\raggedright
Tipo de carga
\end{minipage} & \begin{minipage}[b]{\linewidth}\raggedright
Ejemplos
\end{minipage} \\
\midrule\noalign{}
\endhead
\bottomrule\noalign{}
\endlastfoot
Iluminación & Luminarias interiores, exteriores, decorativas y de
emergencia \\
Tomacorrientes & Equipos domésticos, computadores, cargadores,
electrodomésticos menores \\
Cargas especiales & Cocina eléctrica, ducha eléctrica, calefón, aire
acondicionado, bomba de agua \\
Servicios complementarios & Porteros eléctricos, alarmas, cámaras, redes
de comunicación \\
\end{longtable}
}

En instalaciones civiles, el diseño prioriza la seguridad de las
personas, la continuidad del servicio, la facilidad de uso y el
cumplimiento normativo.

\subsection{Instalaciones eléctricas
industriales}\label{instalaciones-eluxe9ctricas-industriales}

Las instalaciones industriales se aplican en procesos productivos,
talleres, plantas de manufactura, sistemas electromecánicos y entornos
donde existen cargas de potencia, control e instrumentación.

Sus cargas más comunes son:

{\def\LTcaptype{none} % do not increment counter
\begin{longtable}[]{@{}
  >{\raggedright\arraybackslash}p{(\linewidth - 2\tabcolsep) * \real{0.5000}}
  >{\raggedright\arraybackslash}p{(\linewidth - 2\tabcolsep) * \real{0.5000}}@{}}
\toprule\noalign{}
\begin{minipage}[b]{\linewidth}\raggedright
Tipo de carga
\end{minipage} & \begin{minipage}[b]{\linewidth}\raggedright
Ejemplos
\end{minipage} \\
\midrule\noalign{}
\endhead
\bottomrule\noalign{}
\endlastfoot
Cargas de fuerza & Motores, bombas, compresores, ventiladores \\
Cargas de control & Contactores, relés, temporizadores, PLC \\
Cargas de instrumentación & Sensores, transmisores, indicadores \\
Cargas de iluminación industrial & Luminarias de nave, reflectores,
alumbrado de emergencia \\
Sistemas auxiliares & Tableros, bancos de capacitores, variadores de
frecuencia \\
\end{longtable}
}

\subsection{Tabla 1.1. Diferencias entre instalaciones civiles e
industriales}\label{tabla-1.1.-diferencias-entre-instalaciones-civiles-e-industriales}

{\def\LTcaptype{none} % do not increment counter
\begin{longtable}[]{@{}
  >{\raggedright\arraybackslash}p{(\linewidth - 4\tabcolsep) * \real{0.3333}}
  >{\raggedright\arraybackslash}p{(\linewidth - 4\tabcolsep) * \real{0.3333}}
  >{\raggedright\arraybackslash}p{(\linewidth - 4\tabcolsep) * \real{0.3333}}@{}}
\toprule\noalign{}
\begin{minipage}[b]{\linewidth}\raggedright
Criterio
\end{minipage} & \begin{minipage}[b]{\linewidth}\raggedright
Instalaciones civiles
\end{minipage} & \begin{minipage}[b]{\linewidth}\raggedright
Instalaciones industriales
\end{minipage} \\
\midrule\noalign{}
\endhead
\bottomrule\noalign{}
\endlastfoot
Uso principal & Vivienda, comercio y servicios & Producción, control y
operación de maquinaria \\
Tipo de cargas & Iluminación, tomacorrientes y cargas especiales &
Motores, tableros de control, sensores y actuadores \\
Nivel de potencia & Bajo a medio & Medio a alto \\
Complejidad técnica & Moderada & Alta \\
Documentos frecuentes & Plano eléctrico, cuadro de cargas, diagrama
unifilar & Diagramas de potencia, control, mando, PLC e
instrumentación \\
Riesgos principales & Choque eléctrico, sobrecarga, mala puesta a tierra
& Arco eléctrico, fallas de motor, arranque de cargas, fallas de
control \\
Mantenimiento & Correctivo y preventivo básico & Preventivo, predictivo
y correctivo especializado \\
\end{longtable}
}

\begin{tcolorbox}[enhanced jigsaw, arc=.35mm, bottomrule=.15mm, bottomtitle=1mm, breakable, colback=white, colbacktitle=quarto-callout-important-color!10!white, colframe=quarto-callout-important-color-frame, coltitle=black, left=2mm, leftrule=.75mm, opacityback=0, opacitybacktitle=0.6, rightrule=.15mm, title=\textcolor{quarto-callout-important-color}{\faExclamation}\hspace{0.5em}{Importante}, titlerule=0mm, toprule=.15mm, toptitle=1mm]

La principal diferencia entre una instalación civil y una industrial no
es únicamente el tamaño, sino el tipo de carga, el nivel de control, la
criticidad operativa y los riesgos asociados.

\end{tcolorbox}

\begin{center}\rule{0.5\linewidth}{0.5pt}\end{center}

\section{1.2 Importancia de la normalización
eléctrica}\label{importancia-de-la-normalizaciuxf3n-eluxe9ctrica}

\subsection{¿Qué es la normalización
eléctrica?}\label{quuxe9-es-la-normalizaciuxf3n-eluxe9ctrica}

La normalización eléctrica es el conjunto de reglas, criterios,
símbolos, procedimientos y requisitos técnicos que permiten diseñar,
construir, operar e inspeccionar instalaciones eléctricas bajo
parámetros comunes.

Su finalidad es evitar interpretaciones ambiguas y reducir riesgos
técnicos.

\subsection{Importancia técnica}\label{importancia-tuxe9cnica}

La normalización permite:

\begin{enumerate}
\def\labelenumi{\arabic{enumi}.}
\tightlist
\item
  Garantizar condiciones mínimas de seguridad.
\item
  Unificar criterios entre diseñadores, constructores, instaladores y
  fiscalizadores.
\item
  Facilitar la lectura de planos eléctricos.
\item
  Definir características mínimas de materiales y equipos.
\item
  Evitar sobredimensionamientos o subdimensionamientos.
\item
  Reducir el riesgo de incendios de origen eléctrico.
\item
  Facilitar el mantenimiento y ampliación de instalaciones.
\item
  Proteger a usuarios, equipos y edificaciones.
\end{enumerate}

\subsection{Figura 1.2. Función de la normalización
eléctrica}\label{figura-1.2.-funciuxf3n-de-la-normalizaciuxf3n-eluxe9ctrica}

\includegraphics[width=15.8in,height=2.06in]{bloque_01_normalizacion_seguridad_simbologia_files/figure-latex/mermaid-figure-8.png}

\subsection{Ejemplo aplicado}\label{ejemplo-aplicado}

Supóngase que un plano residencial representa un tomacorriente común y
una salida para cocina eléctrica usando el mismo símbolo. El instalador
podría utilizar el mismo conductor y la misma protección en ambos
puntos. Esto generaría un riesgo, porque la cocina eléctrica requiere un
circuito especial, mayor capacidad de corriente y protección
independiente.

La normalización evita este tipo de errores porque obliga a diferenciar
símbolos, circuitos, protecciones y cargas.

\subsection{Tabla 1.2. Problemas frecuentes por falta de
normalización}\label{tabla-1.2.-problemas-frecuentes-por-falta-de-normalizaciuxf3n}

{\def\LTcaptype{none} % do not increment counter
\begin{longtable}[]{@{}
  >{\raggedright\arraybackslash}p{(\linewidth - 4\tabcolsep) * \real{0.3333}}
  >{\raggedright\arraybackslash}p{(\linewidth - 4\tabcolsep) * \real{0.3333}}
  >{\raggedright\arraybackslash}p{(\linewidth - 4\tabcolsep) * \real{0.3333}}@{}}
\toprule\noalign{}
\begin{minipage}[b]{\linewidth}\raggedright
Problema
\end{minipage} & \begin{minipage}[b]{\linewidth}\raggedright
Consecuencia técnica
\end{minipage} & \begin{minipage}[b]{\linewidth}\raggedright
Riesgo asociado
\end{minipage} \\
\midrule\noalign{}
\endhead
\bottomrule\noalign{}
\endlastfoot
Símbolos no normalizados & Mala interpretación del plano & Conexiones
incorrectas \\
Conductores sin código de color & Dificultad para identificar fase,
neutro y tierra & Choque eléctrico \\
Tableros sin rotulación & Mantenimiento inseguro & Desconexión
incorrecta \\
Circuitos sin identificación & Confusión en operación y mantenimiento &
Sobrecarga o mal aislamiento \\
Ausencia de cuadro de cargas & Selección deficiente de protecciones &
Calentamiento o incendio \\
No considerar puesta a tierra & Fallas peligrosas en carcasas metálicas
& Contacto indirecto \\
\end{longtable}
}

\begin{center}\rule{0.5\linewidth}{0.5pt}\end{center}

\section{1.3 Normas técnicas aplicadas a instalaciones
eléctricas}\label{normas-tuxe9cnicas-aplicadas-a-instalaciones-eluxe9ctricas}

\subsection{Normativa principal}\label{normativa-principal}

En Ecuador, una referencia fundamental para instalaciones eléctricas en
edificaciones es la \textbf{NEC-SB-IE Instalaciones Eléctricas},
perteneciente a la Norma Ecuatoriana de la Construcción. Esta norma se
enfoca en instalaciones eléctricas interiores residenciales y establece
criterios mínimos de diseño y seguridad.

La NEC-SB-IE señala que las instalaciones deben contar con protección
frente a:

\begin{itemize}
\tightlist
\item
  choques eléctricos;
\item
  efectos térmicos;
\item
  sobrecorrientes;
\item
  corrientes de falla;
\item
  sobrevoltajes.
\end{itemize}

\subsection{Tabla 1.3. Normas y documentos técnicos de
referencia}\label{tabla-1.3.-normas-y-documentos-tuxe9cnicos-de-referencia}

{\def\LTcaptype{none} % do not increment counter
\begin{longtable}[]{@{}
  >{\raggedright\arraybackslash}p{(\linewidth - 2\tabcolsep) * \real{0.5000}}
  >{\raggedright\arraybackslash}p{(\linewidth - 2\tabcolsep) * \real{0.5000}}@{}}
\toprule\noalign{}
\begin{minipage}[b]{\linewidth}\raggedright
Norma / documento
\end{minipage} & \begin{minipage}[b]{\linewidth}\raggedright
Aplicación en el módulo
\end{minipage} \\
\midrule\noalign{}
\endhead
\bottomrule\noalign{}
\endlastfoot
NEC-SB-IE Instalaciones Eléctricas & Diseño y ejecución de instalaciones
eléctricas interiores residenciales \\
NFPA 70 / National Electrical Code & Criterios de seguridad,
protecciones, conductores y canalizaciones \\
CPE INEN 019 & Código Eléctrico Ecuatoriano \\
IEC 60617 & Simbología gráfica para diagramas eléctricos \\
NTE INEN 2345 & Alambres y cables con aislamiento termoplástico \\
NTE INEN 3098 & Voltajes normalizados \\
IEC 60364 & Instalaciones eléctricas en edificaciones \\
Reglamentos de empresas distribuidoras & Acometidas, medidores y
requisitos de conexión \\
\end{longtable}
}

\begin{tcolorbox}[enhanced jigsaw, arc=.35mm, bottomrule=.15mm, bottomtitle=1mm, breakable, colback=white, colbacktitle=quarto-callout-note-color!10!white, colframe=quarto-callout-note-color-frame, coltitle=black, left=2mm, leftrule=.75mm, opacityback=0, opacitybacktitle=0.6, rightrule=.15mm, title=\textcolor{quarto-callout-note-color}{\faInfo}\hspace{0.5em}{Nota}, titlerule=0mm, toprule=.15mm, toptitle=1mm]

Cuando existan diferencias entre documentos técnicos, se debe revisar la
normativa nacional vigente, los requisitos de la empresa distribuidora y
las especificaciones del proyecto.

\end{tcolorbox}

\subsection{Relación entre norma, plano y
ejecución}\label{relaciuxf3n-entre-norma-plano-y-ejecuciuxf3n}

Una instalación eléctrica segura requiere coherencia entre:

\begin{enumerate}
\def\labelenumi{\arabic{enumi}.}
\tightlist
\item
  normativa aplicable;
\item
  memoria técnica;
\item
  plano eléctrico;
\item
  cuadro de cargas;
\item
  diagrama unifilar;
\item
  selección de materiales;
\item
  ejecución en obra;
\item
  inspección final.
\end{enumerate}

\subsection{Figura 1.3. Flujo técnico de diseño
normativo}\label{figura-1.3.-flujo-tuxe9cnico-de-diseuxf1o-normativo}

\includegraphics[width=18.11in,height=0.98in]{bloque_01_normalizacion_seguridad_simbologia_files/figure-latex/mermaid-figure-7.png}

\begin{center}\rule{0.5\linewidth}{0.5pt}\end{center}

\section{1.4 Simbología eléctrica
normalizada}\label{simbologuxeda-eluxe9ctrica-normalizada}

\subsection{Concepto}\label{concepto}

La simbología eléctrica normalizada permite representar gráficamente los
componentes de una instalación eléctrica. Un símbolo eléctrico resume
información técnica sobre el tipo de dispositivo, su función y su
ubicación dentro del sistema.

Los símbolos se emplean en:

\begin{itemize}
\tightlist
\item
  planos eléctricos en planta;
\item
  diagramas unifilares;
\item
  diagramas multifilares;
\item
  diagramas de control;
\item
  diagramas de potencia;
\item
  planos de tableros;
\item
  esquemas de automatización.
\end{itemize}

\subsection{Tabla 1.4. Simbología básica para instalaciones
civiles}\label{tabla-1.4.-simbologuxeda-buxe1sica-para-instalaciones-civiles}

{\def\LTcaptype{none} % do not increment counter
\begin{longtable}[]{@{}
  >{\raggedright\arraybackslash}p{(\linewidth - 4\tabcolsep) * \real{0.3333}}
  >{\raggedright\arraybackslash}p{(\linewidth - 4\tabcolsep) * \real{0.3333}}
  >{\raggedright\arraybackslash}p{(\linewidth - 4\tabcolsep) * \real{0.3333}}@{}}
\toprule\noalign{}
\begin{minipage}[b]{\linewidth}\raggedright
Elemento
\end{minipage} & \begin{minipage}[b]{\linewidth}\raggedright
Representación textual sugerida
\end{minipage} & \begin{minipage}[b]{\linewidth}\raggedright
Función
\end{minipage} \\
\midrule\noalign{}
\endhead
\bottomrule\noalign{}
\endlastfoot
Punto de luz & PL & Salida para luminaria \\
Interruptor simple & S1 & Control de una luminaria desde un punto \\
Interruptor doble & S2 & Control de dos luminarias o grupos \\
Tomacorriente simple & T1 & Alimentación de carga de uso general \\
Tomacorriente doble & T2 & Alimentación de dos equipos \\
Tomacorriente con tierra & TT & Alimentación con conductor de
protección \\
Salida especial & SE & Alimentación de carga específica \\
Tablero de distribución & TD & Agrupación de protecciones \\
Medidor & M & Registro de energía \\
Puesta a tierra & PT & Conexión de seguridad a tierra \\
\end{longtable}
}

\subsection{Tabla 1.5. Simbología básica para instalaciones
industriales}\label{tabla-1.5.-simbologuxeda-buxe1sica-para-instalaciones-industriales}

{\def\LTcaptype{none} % do not increment counter
\begin{longtable}[]{@{}lll@{}}
\toprule\noalign{}
Elemento & Código frecuente & Función \\
\midrule\noalign{}
\endhead
\bottomrule\noalign{}
\endlastfoot
Motor & M & Convierte energía eléctrica en mecánica \\
Contactor & KM & Maniobra eléctrica de cargas \\
Relé térmico & FR & Protección contra sobrecarga en motores \\
Pulsador de marcha & S1 / START & Orden de encendido \\
Pulsador de paro & S0 / STOP & Orden de apagado \\
Lámpara piloto & H & Señalización visual \\
Final de carrera & SQ & Detección de posición \\
Sensor & B & Detección de variable física \\
PLC & PLC & Control lógico programable \\
Fuente de alimentación & PS & Alimentación de circuitos de control \\
\end{longtable}
}

\subsection{Figura 1.4. Relación entre símbolo, dispositivo y
función}\label{figura-1.4.-relaciuxf3n-entre-suxedmbolo-dispositivo-y-funciuxf3n}

\includegraphics[width=9.19in,height=0.73in]{bloque_01_normalizacion_seguridad_simbologia_files/figure-latex/mermaid-figure-6.png}

\subsection{Criterios de
representación}\label{criterios-de-representaciuxf3n}

Para que la simbología sea técnicamente correcta:

\begin{itemize}
\tightlist
\item
  debe ser legible;
\item
  debe estar incluida en una leyenda;
\item
  debe mantenerse igual en todo el plano;
\item
  debe diferenciar cargas normales y especiales;
\item
  debe identificar circuitos;
\item
  debe relacionarse con el cuadro de cargas;
\item
  debe permitir al instalador ejecutar el trabajo sin ambigüedad.
\end{itemize}

\begin{center}\rule{0.5\linewidth}{0.5pt}\end{center}

\section{1.5 Identificación de símbolos eléctricos en
planos}\label{identificaciuxf3n-de-suxedmbolos-eluxe9ctricos-en-planos}

\subsection{Lectura de planos
eléctricos}\label{lectura-de-planos-eluxe9ctricos}

Un plano eléctrico es un documento técnico que representa la ubicación y
conexión de elementos eléctricos dentro de una edificación.

En instalaciones civiles, el plano eléctrico se suele elaborar sobre la
planta arquitectónica. Sobre ella se ubican luminarias, interruptores,
tomacorrientes, salidas especiales, tablero y canalizaciones.

\subsection{Elementos mínimos de un plano eléctrico
residencial}\label{elementos-muxednimos-de-un-plano-eluxe9ctrico-residencial}

{\def\LTcaptype{none} % do not increment counter
\begin{longtable}[]{@{}
  >{\raggedright\arraybackslash}p{(\linewidth - 4\tabcolsep) * \real{0.3333}}
  >{\raggedright\arraybackslash}p{(\linewidth - 4\tabcolsep) * \real{0.3333}}
  >{\raggedright\arraybackslash}p{(\linewidth - 4\tabcolsep) * \real{0.3333}}@{}}
\toprule\noalign{}
\begin{minipage}[b]{\linewidth}\raggedright
Elemento
\end{minipage} & \begin{minipage}[b]{\linewidth}\raggedright
Descripción
\end{minipage} & \begin{minipage}[b]{\linewidth}\raggedright
Importancia
\end{minipage} \\
\midrule\noalign{}
\endhead
\bottomrule\noalign{}
\endlastfoot
Planta arquitectónica & Base gráfica de la vivienda & Permite ubicar
ambientes y recorridos \\
Puntos de iluminación & Luminarias interiores y exteriores & Define
circuitos de alumbrado \\
Interruptores & Dispositivos de control & Permiten maniobra de
luminarias \\
Tomacorrientes & Salidas de uso general & Alimentan equipos
portátiles \\
Cargas especiales & Cocina, ducha, calefón, bomba & Requieren circuitos
independientes \\
Tablero & Centro de distribución & Aloja protecciones \\
Canalizaciones & Recorrido de conductores & Define ejecución física \\
Cuadro de cargas & Resumen de potencia y corriente & Permite
dimensionamiento \\
Diagrama unifilar & Esquema general del sistema & Muestra jerarquía
eléctrica \\
Leyenda & Símbolos utilizados & Facilita interpretación \\
\end{longtable}
}

\subsection{Figura 1.5. Esquema conceptual de plano eléctrico
residencial}\label{figura-1.5.-esquema-conceptual-de-plano-eluxe9ctrico-residencial}

\includegraphics[width=12.18in,height=5.31in]{bloque_01_normalizacion_seguridad_simbologia_files/figure-latex/mermaid-figure-5.png}

\subsection{Procedimiento recomendado para interpretar un plano
eléctrico}\label{procedimiento-recomendado-para-interpretar-un-plano-eluxe9ctrico}

\begin{enumerate}
\def\labelenumi{\arabic{enumi}.}
\tightlist
\item
  Identificar la planta arquitectónica.
\item
  Localizar el tablero de distribución.
\item
  Reconocer la simbología utilizada.
\item
  Revisar los circuitos de iluminación.
\item
  Revisar los circuitos de tomacorrientes.
\item
  Identificar cargas especiales.
\item
  Verificar recorridos de canalización.
\item
  Revisar el cuadro de cargas.
\item
  Comparar el plano con el diagrama unifilar.
\item
  Confirmar la existencia de sistema de puesta a tierra.
\end{enumerate}

\subsection{Ejemplo de lectura
técnica}\label{ejemplo-de-lectura-tuxe9cnica}

Si en el plano aparece una salida especial en cocina, el estudiante debe
verificar:

\begin{itemize}
\tightlist
\item
  si tiene circuito independiente;
\item
  si aparece en el cuadro de cargas;
\item
  si se asignó potencia adecuada;
\item
  si tiene conductor y protección correspondiente;
\item
  si existe conductor de puesta a tierra;
\item
  si aparece en el diagrama unifilar.
\end{itemize}

\begin{tcolorbox}[enhanced jigsaw, arc=.35mm, bottomrule=.15mm, bottomtitle=1mm, breakable, colback=white, colbacktitle=quarto-callout-tip-color!10!white, colframe=quarto-callout-tip-color-frame, coltitle=black, left=2mm, leftrule=.75mm, opacityback=0, opacitybacktitle=0.6, rightrule=.15mm, title=\textcolor{quarto-callout-tip-color}{\faLightbulb}\hspace{0.5em}{Tip}, titlerule=0mm, toprule=.15mm, toptitle=1mm]

Todo símbolo colocado en planta debe tener respaldo en la leyenda y,
cuando corresponda, en el cuadro de cargas.

\end{tcolorbox}

\begin{center}\rule{0.5\linewidth}{0.5pt}\end{center}

\section{1.6 Normas de seguridad eléctrica y uso de laboratorios
eléctricos}\label{normas-de-seguridad-eluxe9ctrica-y-uso-de-laboratorios-eluxe9ctricos}

\subsection{Seguridad eléctrica}\label{seguridad-eluxe9ctrica}

La seguridad eléctrica consiste en aplicar medidas preventivas para
evitar daños a las personas, equipos e instalaciones. En instalaciones
eléctricas, el riesgo no solo aparece por contacto directo con
conductores energizados, sino también por fallas de aislamiento,
conexiones deficientes, sobrecargas, ausencia de puesta a tierra o
protecciones mal seleccionadas.

\subsection{Seguridad en laboratorio}\label{seguridad-en-laboratorio}

El laboratorio eléctrico es un espacio de aprendizaje experimental. Sin
embargo, al trabajar con tensión, corriente y equipos de medición, se
deben cumplir reglas estrictas.

\subsection{Tabla 1.7. Reglas básicas de seguridad en
laboratorio}\label{tabla-1.7.-reglas-buxe1sicas-de-seguridad-en-laboratorio}

{\def\LTcaptype{none} % do not increment counter
\begin{longtable}[]{@{}
  >{\raggedright\arraybackslash}p{(\linewidth - 2\tabcolsep) * \real{0.5000}}
  >{\raggedright\arraybackslash}p{(\linewidth - 2\tabcolsep) * \real{0.5000}}@{}}
\toprule\noalign{}
\begin{minipage}[b]{\linewidth}\raggedright
Regla
\end{minipage} & \begin{minipage}[b]{\linewidth}\raggedright
Aplicación práctica
\end{minipage} \\
\midrule\noalign{}
\endhead
\bottomrule\noalign{}
\endlastfoot
Revisar el circuito antes de energizar & Verificar conexiones según el
diagrama \\
Confirmar ausencia de tensión & Usar multímetro antes de manipular \\
Seleccionar escala correcta & Evitar daño del instrumento \\
Usar cables en buen estado & Evitar falsos contactos o cortocircuitos \\
No dejar conductores expuestos & Prevenir choque eléctrico \\
Energizar solo con autorización & Mantener control docente \\
Desenergizar antes de modificar & Evitar conexiones bajo tensión \\
Mantener orden en la mesa & Reducir errores de conexión \\
No improvisar puentes & Respetar el esquema de trabajo \\
Reportar fallas de equipos & Prevenir accidentes \\
\end{longtable}
}

\subsection{Figura 1.6. Secuencia segura de trabajo en
laboratorio}\label{figura-1.6.-secuencia-segura-de-trabajo-en-laboratorio}

\includegraphics[width=23.4in,height=0.98in]{bloque_01_normalizacion_seguridad_simbologia_files/figure-latex/mermaid-figure-4.png}

\subsection{Uso seguro del
multímetro}\label{uso-seguro-del-multuxedmetro}

El multímetro es uno de los instrumentos más usados en instalaciones
eléctricas. Permite medir voltaje, corriente, resistencia y continuidad,
según el modelo.

\subsection{Tabla 1.8. Uso básico del
multímetro}\label{tabla-1.8.-uso-buxe1sico-del-multuxedmetro}

{\def\LTcaptype{none} % do not increment counter
\begin{longtable}[]{@{}
  >{\raggedright\arraybackslash}p{(\linewidth - 6\tabcolsep) * \real{0.2500}}
  >{\raggedright\arraybackslash}p{(\linewidth - 6\tabcolsep) * \real{0.2500}}
  >{\raggedright\arraybackslash}p{(\linewidth - 6\tabcolsep) * \real{0.2500}}
  >{\raggedright\arraybackslash}p{(\linewidth - 6\tabcolsep) * \real{0.2500}}@{}}
\toprule\noalign{}
\begin{minipage}[b]{\linewidth}\raggedright
Magnitud
\end{minipage} & \begin{minipage}[b]{\linewidth}\raggedright
Posición del instrumento
\end{minipage} & \begin{minipage}[b]{\linewidth}\raggedright
Forma de conexión
\end{minipage} & \begin{minipage}[b]{\linewidth}\raggedright
Condición de seguridad
\end{minipage} \\
\midrule\noalign{}
\endhead
\bottomrule\noalign{}
\endlastfoot
Voltaje & V AC / V DC & En paralelo & Circuito energizado con cuidado \\
Corriente & A / mA & En serie & Revisar escala y borne correcto \\
Resistencia & Ω & En paralelo al elemento & Circuito desenergizado \\
Continuidad & Símbolo de buzzer & Entre dos puntos & Circuito
desenergizado \\
\end{longtable}
}

\begin{tcolorbox}[enhanced jigsaw, arc=.35mm, bottomrule=.15mm, bottomtitle=1mm, breakable, colback=white, colbacktitle=quarto-callout-warning-color!10!white, colframe=quarto-callout-warning-color-frame, coltitle=black, left=2mm, leftrule=.75mm, opacityback=0, opacitybacktitle=0.6, rightrule=.15mm, title=\textcolor{quarto-callout-warning-color}{\faExclamationTriangle}\hspace{0.5em}{Advertencia}, titlerule=0mm, toprule=.15mm, toptitle=1mm]

Nunca se debe medir resistencia o continuidad en un circuito energizado.
Esta acción puede dañar el instrumento y poner en riesgo al usuario.

\end{tcolorbox}

\subsection{Equipo de protección
personal}\label{equipo-de-protecciuxf3n-personal}

{\def\LTcaptype{none} % do not increment counter
\begin{longtable}[]{@{}ll@{}}
\toprule\noalign{}
Elemento & Función \\
\midrule\noalign{}
\endhead
\bottomrule\noalign{}
\endlastfoot
Gafas de seguridad & Protección ante chispas o partículas \\
Guantes aislantes & Protección frente a contacto eléctrico \\
Calzado dieléctrico & Aislamiento respecto al piso \\
Herramientas aisladas & Reducción de riesgo por contacto accidental \\
Ropa ajustada & Evita enganches o contacto accidental \\
Señalización & Advierte presencia de trabajo eléctrico \\
\end{longtable}
}

\begin{center}\rule{0.5\linewidth}{0.5pt}\end{center}

\section{1.7 Riesgos eléctricos y medidas de
prevención}\label{riesgos-eluxe9ctricos-y-medidas-de-prevenciuxf3n}

\subsection{Concepto de riesgo
eléctrico}\label{concepto-de-riesgo-eluxe9ctrico}

Un riesgo eléctrico es la posibilidad de que una persona, equipo o
instalación sufra daño debido a la presencia de energía eléctrica en
condiciones inseguras.

Los riesgos eléctricos pueden presentarse por:

\begin{itemize}
\tightlist
\item
  contacto directo;
\item
  contacto indirecto;
\item
  sobrecarga;
\item
  cortocircuito;
\item
  falla a tierra;
\item
  arco eléctrico;
\item
  conexiones deficientes;
\item
  ausencia de protecciones;
\item
  mala selección de conductores;
\item
  canalizaciones inadecuadas.
\end{itemize}

\subsection{Tabla 1.9. Tipos de riesgos
eléctricos}\label{tabla-1.9.-tipos-de-riesgos-eluxe9ctricos}

{\def\LTcaptype{none} % do not increment counter
\begin{longtable}[]{@{}
  >{\raggedright\arraybackslash}p{(\linewidth - 6\tabcolsep) * \real{0.2500}}
  >{\raggedright\arraybackslash}p{(\linewidth - 6\tabcolsep) * \real{0.2500}}
  >{\raggedright\arraybackslash}p{(\linewidth - 6\tabcolsep) * \real{0.2500}}
  >{\raggedright\arraybackslash}p{(\linewidth - 6\tabcolsep) * \real{0.2500}}@{}}
\toprule\noalign{}
\begin{minipage}[b]{\linewidth}\raggedright
Riesgo
\end{minipage} & \begin{minipage}[b]{\linewidth}\raggedright
Causa frecuente
\end{minipage} & \begin{minipage}[b]{\linewidth}\raggedright
Consecuencia
\end{minipage} & \begin{minipage}[b]{\linewidth}\raggedright
Prevención
\end{minipage} \\
\midrule\noalign{}
\endhead
\bottomrule\noalign{}
\endlastfoot
Choque eléctrico & Contacto con parte energizada & Lesión o muerte &
Aislamiento, señalización, EPP \\
Contacto indirecto & Carcasa energizada por falla & Electrocución &
Puesta a tierra, protección diferencial \\
Sobrecarga & Exceso de corriente en conductor & Calentamiento & Cálculo
de carga y protección adecuada \\
Cortocircuito & Unión accidental entre conductores & Corriente elevada &
Breaker, fusible, aislamiento correcto \\
Falla a tierra & Corriente hacia tierra no prevista & Riesgo en equipos
metálicos & Sistema de puesta a tierra \\
Arco eléctrico & Falla de alta energía & Quemaduras, daño de equipos &
Mantenimiento, EPP, protecciones \\
Incendio eléctrico & Calor por mala conexión & Daño material &
Inspección y conexiones firmes \\
\end{longtable}
}

\subsection{Figura 1.7. Relación entre falla eléctrica y
protección}\label{figura-1.7.-relaciuxf3n-entre-falla-eluxe9ctrica-y-protecciuxf3n}

\includegraphics[width=6.96in,height=5.31in]{bloque_01_normalizacion_seguridad_simbologia_files/figure-latex/mermaid-figure-3.png}

\subsection{Sobrecarga}\label{sobrecarga}

La sobrecarga ocurre cuando un circuito conduce una corriente mayor a la
prevista durante un tiempo prolongado. Esto puede suceder por conectar
demasiadas cargas a un mismo circuito o por seleccionar conductores de
calibre insuficiente.

Ejemplo:

Si un circuito de tomacorrientes fue diseñado para 20 A y se conectan
equipos que demandan 28 A, el conductor puede calentarse excesivamente
si la protección no actúa adecuadamente.

\subsection{Cortocircuito}\label{cortocircuito}

El cortocircuito es una falla de baja impedancia entre conductores
activos o entre fase y tierra. Genera una corriente muy elevada en un
tiempo muy corto.

Sus causas frecuentes son:

\begin{itemize}
\tightlist
\item
  aislamiento deteriorado;
\item
  conexiones flojas;
\item
  humedad;
\item
  daño mecánico en conductores;
\item
  errores de conexión;
\item
  contacto accidental entre fase y neutro.
\end{itemize}

\subsection{Falla a tierra}\label{falla-a-tierra}

La falla a tierra ocurre cuando una corriente no deseada fluye hacia
elementos metálicos o hacia tierra. Es especialmente peligrosa cuando
energiza carcasas de equipos.

La puesta a tierra proporciona un camino seguro para esa corriente de
falla y permite la operación de las protecciones.

\subsection{Medidas generales de
prevención}\label{medidas-generales-de-prevenciuxf3n}

{\def\LTcaptype{none} % do not increment counter
\begin{longtable}[]{@{}
  >{\raggedright\arraybackslash}p{(\linewidth - 2\tabcolsep) * \real{0.5000}}
  >{\raggedright\arraybackslash}p{(\linewidth - 2\tabcolsep) * \real{0.5000}}@{}}
\toprule\noalign{}
\begin{minipage}[b]{\linewidth}\raggedright
Medida
\end{minipage} & \begin{minipage}[b]{\linewidth}\raggedright
Aplicación
\end{minipage} \\
\midrule\noalign{}
\endhead
\bottomrule\noalign{}
\endlastfoot
Dimensionar correctamente conductores & Evita calentamiento \\
Seleccionar protecciones adecuadas & Permite desconexión ante fallas \\
Usar puesta a tierra & Reduce riesgo de contacto indirecto \\
Usar protección diferencial & Protege personas ante fugas de
corriente \\
Aislar conexiones & Evita contactos accidentales \\
Rotular tableros & Facilita mantenimiento seguro \\
Separar circuitos & Evita sobrecarga y facilita operación \\
Realizar inspecciones & Detecta fallas antes de accidentes \\
\end{longtable}
}

\begin{center}\rule{0.5\linewidth}{0.5pt}\end{center}

\section{Integración del Bloque 1 con la práctica
experimental}\label{integraciuxf3n-del-bloque-1-con-la-pruxe1ctica-experimental}

\subsection{Práctica del bloque}\label{pruxe1ctica-del-bloque}

La práctica oficial del bloque consiste en el diseño de un plano
eléctrico residencial. Esta actividad integra normalización, simbología,
seguridad y criterios iniciales de diseño.

\subsection{Figura 1.8. Flujo de trabajo de la práctica
residencial}\label{figura-1.8.-flujo-de-trabajo-de-la-pruxe1ctica-residencial}

\includegraphics[width=2.88in,height=9.9in]{bloque_01_normalizacion_seguridad_simbologia_files/figure-latex/mermaid-figure-2.png}

\subsection{Entregables de la
práctica}\label{entregables-de-la-pruxe1ctica}

{\def\LTcaptype{none} % do not increment counter
\begin{longtable}[]{@{}
  >{\raggedright\arraybackslash}p{(\linewidth - 2\tabcolsep) * \real{0.5000}}
  >{\raggedright\arraybackslash}p{(\linewidth - 2\tabcolsep) * \real{0.5000}}@{}}
\toprule\noalign{}
\begin{minipage}[b]{\linewidth}\raggedright
Entregable
\end{minipage} & \begin{minipage}[b]{\linewidth}\raggedright
Contenido mínimo
\end{minipage} \\
\midrule\noalign{}
\endhead
\bottomrule\noalign{}
\endlastfoot
Plano eléctrico residencial & Luminarias, interruptores, tomacorrientes,
cargas especiales y tablero \\
Cuadro de cargas & Circuito, descripción, potencia, corriente, conductor
y protección \\
Diagrama unifilar & Medidor, tablero, protecciones y circuitos
derivados \\
Cuadro de simbología & Símbolos utilizados en el plano \\
Memoria técnica & Criterios normativos, cálculos y justificación \\
Conclusiones & Análisis técnico del diseño realizado \\
\end{longtable}
}

\subsection{Cuadro de cargas modelo}\label{cuadro-de-cargas-modelo}

{\def\LTcaptype{none} % do not increment counter
\begin{longtable}[]{@{}
  >{\raggedright\arraybackslash}p{(\linewidth - 16\tabcolsep) * \real{0.0938}}
  >{\raggedright\arraybackslash}p{(\linewidth - 16\tabcolsep) * \real{0.0938}}
  >{\raggedleft\arraybackslash}p{(\linewidth - 16\tabcolsep) * \real{0.1250}}
  >{\raggedleft\arraybackslash}p{(\linewidth - 16\tabcolsep) * \real{0.1250}}
  >{\raggedleft\arraybackslash}p{(\linewidth - 16\tabcolsep) * \real{0.1250}}
  >{\raggedleft\arraybackslash}p{(\linewidth - 16\tabcolsep) * \real{0.1250}}
  >{\raggedleft\arraybackslash}p{(\linewidth - 16\tabcolsep) * \real{0.1250}}
  >{\raggedright\arraybackslash}p{(\linewidth - 16\tabcolsep) * \real{0.0938}}
  >{\raggedright\arraybackslash}p{(\linewidth - 16\tabcolsep) * \real{0.0938}}@{}}
\toprule\noalign{}
\begin{minipage}[b]{\linewidth}\raggedright
Circuito
\end{minipage} & \begin{minipage}[b]{\linewidth}\raggedright
Descripción
\end{minipage} & \begin{minipage}[b]{\linewidth}\raggedleft
Cantidad
\end{minipage} & \begin{minipage}[b]{\linewidth}\raggedleft
Potencia unitaria (W)
\end{minipage} & \begin{minipage}[b]{\linewidth}\raggedleft
Potencia total (W)
\end{minipage} & \begin{minipage}[b]{\linewidth}\raggedleft
Voltaje (V)
\end{minipage} & \begin{minipage}[b]{\linewidth}\raggedleft
Corriente (A)
\end{minipage} & \begin{minipage}[b]{\linewidth}\raggedright
Conductor sugerido
\end{minipage} & \begin{minipage}[b]{\linewidth}\raggedright
Protección
\end{minipage} \\
\midrule\noalign{}
\endhead
\bottomrule\noalign{}
\endlastfoot
C1 & Iluminación planta baja & 8 & 100 & 800 & 120 & 6.67 & 14 AWG & 15
A \\
C2 & Iluminación planta alta & 6 & 100 & 600 & 120 & 5.00 & 14 AWG & 15
A \\
C3 & Tomacorrientes generales & 8 & 200 & 1600 & 120 & 13.33 & 12 AWG &
20 A \\
C4 & Tomacorrientes cocina & 4 & 200 & 800 & 120 & 6.67 & 12 AWG & 20
A \\
C5 & Cocina eléctrica & 1 & 6000 & 6000 & 240 & 25.00 & Según cálculo &
Bipolar \\
C6 & Ducha eléctrica & 1 & 3500 & 3500 & 240 & 14.58 & Según cálculo &
Bipolar \\
\end{longtable}
}

\begin{tcolorbox}[enhanced jigsaw, arc=.35mm, bottomrule=.15mm, bottomtitle=1mm, breakable, colback=white, colbacktitle=quarto-callout-important-color!10!white, colframe=quarto-callout-important-color-frame, coltitle=black, left=2mm, leftrule=.75mm, opacityback=0, opacitybacktitle=0.6, rightrule=.15mm, title=\textcolor{quarto-callout-important-color}{\faExclamation}\hspace{0.5em}{Importante}, titlerule=0mm, toprule=.15mm, toptitle=1mm]

Los valores del cuadro son referenciales para fines académicos. En un
proyecto real deben verificarse con la normativa vigente, condiciones de
instalación, temperatura, caída de tensión, tipo de aislamiento,
longitud del circuito y especificaciones del fabricante.

\end{tcolorbox}

\begin{center}\rule{0.5\linewidth}{0.5pt}\end{center}

\section{Actividad en clase}\label{actividad-en-clase}

\subsection{Análisis de una vivienda
básica}\label{anuxe1lisis-de-una-vivienda-buxe1sica}

A partir de un plano arquitectónico asignado por el docente, el
estudiante deberá:

\begin{enumerate}
\def\labelenumi{\arabic{enumi}.}
\tightlist
\item
  Identificar ambientes.
\item
  Proponer puntos de iluminación.
\item
  Ubicar interruptores.
\item
  Ubicar tomacorrientes.
\item
  Identificar cargas especiales.
\item
  Proponer ubicación del tablero.
\item
  Elaborar una leyenda básica de símbolos.
\item
  Identificar riesgos eléctricos.
\item
  Proponer medidas de prevención.
\item
  Justificar el criterio de separación de circuitos.
\end{enumerate}

\subsection{Producto esperado}\label{producto-esperado}

Un esquema preliminar de instalación eléctrica residencial con
simbología básica, identificación de circuitos y observaciones de
seguridad.

\begin{center}\rule{0.5\linewidth}{0.5pt}\end{center}

\section{Actividad autónoma}\label{actividad-autuxf3noma}

\subsection{Ficha de revisión
normativa}\label{ficha-de-revisiuxf3n-normativa}

El estudiante revisará la NEC-SB-IE Instalaciones Eléctricas y elaborará
una ficha técnica con:

\begin{enumerate}
\def\labelenumi{\arabic{enumi}.}
\tightlist
\item
  objetivo de la norma;
\item
  campo de aplicación;
\item
  normas referenciadas;
\item
  criterios de circuitos;
\item
  criterios de conductores;
\item
  criterios de protecciones;
\item
  criterios de puesta a tierra;
\item
  criterios de ubicación de interruptores y tomacorrientes;
\item
  tres conclusiones técnicas.
\end{enumerate}

Extensión sugerida: 2 a 3 páginas.

\begin{center}\rule{0.5\linewidth}{0.5pt}\end{center}

\section{Preguntas de repaso}\label{preguntas-de-repaso}

\begin{enumerate}
\def\labelenumi{\arabic{enumi}.}
\tightlist
\item
  ¿Qué es una instalación eléctrica?
\item
  ¿Qué diferencia existe entre una instalación civil y una industrial?
\item
  ¿Por qué es importante la normalización eléctrica?
\item
  ¿Qué normas se pueden aplicar en instalaciones eléctricas?
\item
  ¿Qué función cumple la simbología eléctrica?
\item
  ¿Qué elementos debe contener un plano eléctrico residencial?
\item
  ¿Qué es un cuadro de cargas?
\item
  ¿Qué es un diagrama unifilar?
\item
  ¿Qué diferencia existe entre sobrecarga y cortocircuito?
\item
  ¿Por qué la puesta a tierra es una medida de seguridad?
\item
  ¿Qué riesgos existen al trabajar en laboratorio eléctrico?
\item
  ¿Por qué no se debe medir resistencia en un circuito energizado?
\item
  ¿Qué información debe contener una leyenda de simbología?
\item
  ¿Por qué se deben separar circuitos de iluminación, tomacorrientes y
  cargas especiales?
\item
  ¿Qué relación existe entre carga, conductor y protección?
\end{enumerate}

\begin{center}\rule{0.5\linewidth}{0.5pt}\end{center}

\section{Cierre del bloque}\label{cierre-del-bloque}

El Bloque 1 permite al estudiante comprender que una instalación
eléctrica segura depende de la correcta aplicación de normas,
simbología, criterios de diseño y medidas de prevención. Estos
fundamentos son necesarios para avanzar hacia el análisis de diagramas
eléctricos, mediciones, instrumentación y automatización industrial.

En los siguientes bloques se profundizará en la interpretación de
diagramas unifilares, multifilares, diagramas de control y potencia, así
como en el uso de instrumentos de medición y sistemas básicos de
automatización.

\bookmarksetup{startatroot}

\chapter{Bloque 2: Diagramas Eléctricos e
Instrumentación}\label{bloque-2-diagramas-eluxe9ctricos-e-instrumentaciuxf3n}

\section{Resultado de aprendizaje del
bloque}\label{resultado-de-aprendizaje-del-bloque-1}

El estudiante interpreta diagramas y planos eléctricos, realiza
mediciones de magnitudes eléctricas, utiliza instrumentos adecuados y
analiza exactitud, precisión y errores de medición para evaluar sistemas
eléctricos con criterio técnico.

\section{Objetivo específico}\label{objetivo-especuxedfico}

Analizar e interpretar diagramas eléctricos de instalaciones civiles e
industriales, aplicando técnicas de medición eléctrica y uso de
instrumentos para evaluar la exactitud, precisión y errores en los
resultados obtenidos.

\section{2.1 Tipos de diagramas
eléctricos}\label{tipos-de-diagramas-eluxe9ctricos}

Los diagramas eléctricos representan la estructura, funcionamiento o
conexión de un sistema. Entre los más utilizados se encuentran:

\begin{itemize}
\tightlist
\item
  Diagrama unifilar.
\item
  Diagrama multifilar.
\item
  Diagrama de control.
\item
  Diagrama de potencia.
\item
  Plano eléctrico arquitectónico.
\item
  Diagrama de conexión.
\item
  Esquema de tablero.
\end{itemize}

\section{2.2 Diagramas unifilares y
multifilares}\label{diagramas-unifilares-y-multifilares}

El \textbf{diagrama unifilar} representa un sistema eléctrico mediante
una sola línea por circuito o alimentador. Es útil para mostrar
acometida, medidor, protección principal, tableros, circuitos derivados
y cargas.

El \textbf{diagrama multifilar} representa cada conductor de manera
individual. Permite visualizar fases, neutro, tierra, retornos,
contactos y conexiones específicas.

{\def\LTcaptype{none} % do not increment counter
\begin{longtable}[]{@{}
  >{\raggedright\arraybackslash}p{(\linewidth - 4\tabcolsep) * \real{0.3333}}
  >{\raggedright\arraybackslash}p{(\linewidth - 4\tabcolsep) * \real{0.3333}}
  >{\raggedright\arraybackslash}p{(\linewidth - 4\tabcolsep) * \real{0.3333}}@{}}
\toprule\noalign{}
\begin{minipage}[b]{\linewidth}\raggedright
Criterio
\end{minipage} & \begin{minipage}[b]{\linewidth}\raggedright
Diagrama unifilar
\end{minipage} & \begin{minipage}[b]{\linewidth}\raggedright
Diagrama multifilar
\end{minipage} \\
\midrule\noalign{}
\endhead
\bottomrule\noalign{}
\endlastfoot
Nivel de detalle & General & Alto \\
Uso principal & Diseño, memoria técnica, tableros & Cableado y
montaje \\
Representación de conductores & Simplificada & Individual \\
Aplicación & Distribución eléctrica & Conexión real de circuitos \\
\end{longtable}
}

\section{2.3 Diagramas de control y
potencia}\label{diagramas-de-control-y-potencia}

En instalaciones industriales, el circuito de potencia alimenta cargas
como motores, resistencias, bombas o compresores. El circuito de control
gobierna la operación mediante pulsadores, sensores, relés, contactores,
temporizadores o PLC.

\section{2.4 Interpretación de planos
eléctricos}\label{interpretaciuxf3n-de-planos-eluxe9ctricos}

Para interpretar planos eléctricos se debe revisar:

\begin{enumerate}
\def\labelenumi{\arabic{enumi}.}
\tightlist
\item
  Escala y orientación del plano.
\item
  Cuadro de simbología.
\item
  Ubicación de salidas eléctricas.
\item
  Identificación de circuitos.
\item
  Cuadro de cargas.
\item
  Diagrama unifilar.
\item
  Notas técnicas y especificaciones.
\item
  Criterios de seguridad y puesta a tierra.
\end{enumerate}

\section{2.5 Introducción a las mediciones
eléctricas}\label{introducciuxf3n-a-las-mediciones-eluxe9ctricas}

La medición eléctrica permite verificar el comportamiento real de un
sistema. Las magnitudes básicas son voltaje, corriente, resistencia,
potencia y energía.

\section{2.6 Magnitudes eléctricas fundamentales e instrumentos de
medición}\label{magnitudes-eluxe9ctricas-fundamentales-e-instrumentos-de-mediciuxf3n}

{\def\LTcaptype{none} % do not increment counter
\begin{longtable}[]{@{}lllll@{}}
\toprule\noalign{}
Magnitud & Símbolo & Unidad & Instrumento & Conexión básica \\
\midrule\noalign{}
\endhead
\bottomrule\noalign{}
\endlastfoot
Voltaje & V & Voltio & Voltímetro & Paralelo \\
Corriente & I & Amperio & Amperímetro & Serie \\
Resistencia & R & Ohmio & Óhmetro & Circuito desenergizado \\
Potencia & P & Watt & Vatímetro / multímetro & Según instrumento \\
Energía & W o E & kWh & Medidor de energía & Instalación fija \\
\end{longtable}
}

\begin{tcolorbox}[enhanced jigsaw, arc=.35mm, bottomrule=.15mm, bottomtitle=1mm, breakable, colback=white, colbacktitle=quarto-callout-warning-color!10!white, colframe=quarto-callout-warning-color-frame, coltitle=black, left=2mm, leftrule=.75mm, opacityback=0, opacitybacktitle=0.6, rightrule=.15mm, title=\textcolor{quarto-callout-warning-color}{\faExclamationTriangle}\hspace{0.5em}{Seguridad de medición}, titlerule=0mm, toprule=.15mm, toptitle=1mm]

Nunca se debe medir resistencia en un circuito energizado. Primero se
verifica ausencia de tensión y condiciones seguras de trabajo.

\end{tcolorbox}

\section{2.7 Procedimientos de
medición}\label{procedimientos-de-mediciuxf3n}

Procedimiento general:

\begin{enumerate}
\def\labelenumi{\arabic{enumi}.}
\tightlist
\item
  Identificar la magnitud a medir.
\item
  Seleccionar el instrumento adecuado.
\item
  Verificar categoría, rango y estado del instrumento.
\item
  Conectar de acuerdo con la magnitud.
\item
  Registrar lectura y condiciones de medición.
\item
  Interpretar el resultado.
\item
  Desenergizar y ordenar el área de trabajo.
\end{enumerate}

\section{2.8 Exactitud y precisión en las
mediciones}\label{exactitud-y-precisiuxf3n-en-las-mediciones}

La \textbf{exactitud} indica qué tan cerca está una medición del valor
real o aceptado. La \textbf{precisión} indica qué tan repetibles son las
mediciones entre sí.

{\def\LTcaptype{none} % do not increment counter
\begin{longtable}[]{@{}
  >{\raggedright\arraybackslash}p{(\linewidth - 2\tabcolsep) * \real{0.5000}}
  >{\raggedright\arraybackslash}p{(\linewidth - 2\tabcolsep) * \real{0.5000}}@{}}
\toprule\noalign{}
\begin{minipage}[b]{\linewidth}\raggedright
Caso
\end{minipage} & \begin{minipage}[b]{\linewidth}\raggedright
Interpretación
\end{minipage} \\
\midrule\noalign{}
\endhead
\bottomrule\noalign{}
\endlastfoot
Alta exactitud y alta precisión & Mediciones confiables \\
Alta precisión y baja exactitud & Error sistemático probable \\
Baja precisión & Lecturas dispersas \\
Baja exactitud y baja precisión & Procedimiento o instrumento
deficiente \\
\end{longtable}
}

\section{2.9 Errores en las mediciones}\label{errores-en-las-mediciones}

Los errores pueden clasificarse como:

\begin{itemize}
\tightlist
\item
  Sistemáticos: asociados al instrumento, calibración o método.
\item
  Aleatorios: variaciones no controladas durante la medición.
\item
  Humanos: lectura incorrecta, mala conexión o selección inadecuada de
  escala.
\item
  Ambientales: temperatura, humedad, interferencias o condiciones del
  entorno.
\end{itemize}

\subsection{Cálculo básico del error
absoluto}\label{cuxe1lculo-buxe1sico-del-error-absoluto}

\[
E_a = |X_m - X_v|
\]

Donde:

\begin{itemize}
\tightlist
\item
  \(E_a\) = error absoluto.
\item
  \(X_m\) = valor medido.
\item
  \(X_v\) = valor verdadero o de referencia.
\end{itemize}

\subsection{Cálculo básico del error
porcentual}\label{cuxe1lculo-buxe1sico-del-error-porcentual}

\[
E_{\%} = \frac{|X_m - X_v|}{X_v} \times 100
\]

\section{Actividad de aula}\label{actividad-de-aula}

Realice tres mediciones de voltaje en una fuente de laboratorio y
calcule:

\begin{enumerate}
\def\labelenumi{\arabic{enumi}.}
\tightlist
\item
  Promedio.
\item
  Error absoluto.
\item
  Error porcentual.
\item
  Interpretación de exactitud y precisión.
\end{enumerate}

\section{Práctica vinculada}\label{pruxe1ctica-vinculada}

\textbf{Práctica 2:} Diseño de plano eléctrico de una instalación de
tipo industrial.

\bookmarksetup{startatroot}

\chapter{Bloque 3: Automatización}\label{bloque-3-automatizaciuxf3n}

\section{Resultado de aprendizaje del
bloque}\label{resultado-de-aprendizaje-del-bloque-2}

El estudiante analiza e implementa soluciones básicas de automatización
en sistemas eléctricos mediante principios de control industrial,
identificación de arquitectura y componentes de un PLC y aplicaciones en
procesos industriales simples.

\section{Objetivo específico}\label{objetivo-especuxedfico-1}

Desarrollar competencias para comprender y aplicar principios de
automatización industrial y sistemas de control mediante el estudio de
controladores lógicos programables, arquitectura, componentes y
aplicaciones básicas en sistemas eléctricos.

\section{3.1 Introducción a la automatización
industrial}\label{introducciuxf3n-a-la-automatizaciuxf3n-industrial}

La automatización industrial consiste en el uso de dispositivos de
control, sensores, actuadores, interfaces y programas para ejecutar
procesos con mínima intervención humana. Su aplicación mejora
productividad, seguridad, repetibilidad, control de calidad y
continuidad de operación.

\section{3.2 Sistemas de control en procesos
industriales}\label{sistemas-de-control-en-procesos-industriales}

Un sistema de control integra entradas, procesamiento y salidas.

{\def\LTcaptype{none} % do not increment counter
\begin{longtable}[]{@{}ll@{}}
\toprule\noalign{}
Elemento & Función \\
\midrule\noalign{}
\endhead
\bottomrule\noalign{}
\endlastfoot
Sensor & Detecta una condición física o eléctrica \\
Pulsador & Permite una orden manual \\
Controlador & Procesa señales y ejecuta lógica \\
Actuador & Ejecuta una acción física \\
Contactor & Maniobra cargas de potencia \\
Protección & Desconecta ante fallas \\
\end{longtable}
}

\section{3.3 Introducción a los controladores lógicos
programables}\label{introducciuxf3n-a-los-controladores-luxf3gicos-programables}

Un PLC es un dispositivo electrónico diseñado para controlar procesos
mediante lógica programada. Sustituye parcial o totalmente la lógica
cableada de relés, permitiendo modificar secuencias sin cambiar todo el
cableado físico.

\section{3.4 Componentes y arquitectura básica de un
PLC}\label{componentes-y-arquitectura-buxe1sica-de-un-plc}

Un PLC básico incluye:

\begin{itemize}
\tightlist
\item
  Fuente de alimentación.
\item
  CPU.
\item
  Memoria.
\item
  Entradas digitales o analógicas.
\item
  Salidas digitales o analógicas.
\item
  Módulos de comunicación.
\item
  Interfaz de programación.
\end{itemize}

\section{3.5 Aplicaciones básicas de PLC en sistemas
eléctricos}\label{aplicaciones-buxe1sicas-de-plc-en-sistemas-eluxe9ctricos}

Ejemplos de aplicación:

\begin{itemize}
\tightlist
\item
  Arranque y paro de motores.
\item
  Control de iluminación.
\item
  Temporización de cargas.
\item
  Control de bombas.
\item
  Secuencias automáticas.
\item
  Señalización de estados.
\item
  Transferencias automáticas simples.
\end{itemize}

\section{Temporizadores básicos}\label{temporizadores-buxe1sicos}

{\def\LTcaptype{none} % do not increment counter
\begin{longtable}[]{@{}
  >{\raggedright\arraybackslash}p{(\linewidth - 4\tabcolsep) * \real{0.3333}}
  >{\raggedright\arraybackslash}p{(\linewidth - 4\tabcolsep) * \real{0.3333}}
  >{\raggedright\arraybackslash}p{(\linewidth - 4\tabcolsep) * \real{0.3333}}@{}}
\toprule\noalign{}
\begin{minipage}[b]{\linewidth}\raggedright
Temporizador
\end{minipage} & \begin{minipage}[b]{\linewidth}\raggedright
Descripción
\end{minipage} & \begin{minipage}[b]{\linewidth}\raggedright
Aplicación típica
\end{minipage} \\
\midrule\noalign{}
\endhead
\bottomrule\noalign{}
\endlastfoot
TON & Retardo a la conexión & Encendido diferido de una carga \\
TOF & Retardo a la desconexión & Apagado diferido \\
Pulso & Genera activación por tiempo definido & Señalización o eventos
breves \\
\end{longtable}
}

\section{Ejemplo conceptual: retardo a la
conexión}\label{ejemplo-conceptual-retardo-a-la-conexiuxf3n}

Una lámpara piloto debe encenderse cinco segundos después de presionar
un pulsador de marcha. La lógica básica es:

\begin{enumerate}
\def\labelenumi{\arabic{enumi}.}
\tightlist
\item
  Entrada I1 activa el temporizador TON.
\item
  El temporizador cuenta 5 segundos.
\item
  Al cumplirse el tiempo, se activa Q1.
\item
  Un pulsador de paro reinicia el proceso.
\end{enumerate}

\section{Actividad de aula}\label{actividad-de-aula-1}

Diseñe la lógica básica para:

\begin{enumerate}
\def\labelenumi{\arabic{enumi}.}
\tightlist
\item
  Encendido de una lámpara con retardo de 5 segundos.
\item
  Apagado con pulsador de paro.
\item
  Señalización de estado activo.
\item
  Representación del diagrama de conexión.
\end{enumerate}

\section{Práctica vinculada}\label{pruxe1ctica-vinculada-1}

\textbf{Práctica 3:} Fundamentos de programación con PLCs usando Siemens
LOGO!.

\bookmarksetup{startatroot}

\chapter{Prácticas de laboratorio}\label{pruxe1cticas-de-laboratorio}

\section{Organización general}\label{organizaciuxf3n-general}

Las prácticas están alineadas con los tres bloques de la asignatura y
con el componente práctico-experimental.

{\def\LTcaptype{none} % do not increment counter
\begin{longtable}[]{@{}
  >{\raggedright\arraybackslash}p{(\linewidth - 6\tabcolsep) * \real{0.2500}}
  >{\raggedright\arraybackslash}p{(\linewidth - 6\tabcolsep) * \real{0.2500}}
  >{\raggedright\arraybackslash}p{(\linewidth - 6\tabcolsep) * \real{0.2500}}
  >{\raggedright\arraybackslash}p{(\linewidth - 6\tabcolsep) * \real{0.2500}}@{}}
\toprule\noalign{}
\begin{minipage}[b]{\linewidth}\raggedright
Práctica
\end{minipage} & \begin{minipage}[b]{\linewidth}\raggedright
Bloque
\end{minipage} & \begin{minipage}[b]{\linewidth}\raggedright
Título
\end{minipage} & \begin{minipage}[b]{\linewidth}\raggedright
Producto esperado
\end{minipage} \\
\midrule\noalign{}
\endhead
\bottomrule\noalign{}
\endlastfoot
1 & Bloque 1 & Diseño de plano eléctrico residencial & Plano, cuadro de
cargas, diagrama unifilar e informe \\
2 & Bloque 2 & Diseño de plano eléctrico industrial & Plano industrial,
diagramas de potencia/control e informe \\
3 & Bloque 3 & Fundamentos de programación con PLCs & Programa básico,
conexión física e informe \\
\end{longtable}
}

\section{Práctica 1: Diseño de plano eléctrico
residencial}\label{pruxe1ctica-1-diseuxf1o-de-plano-eluxe9ctrico-residencial}

\subsection{Objetivo}\label{objetivo}

Diseñar el plano eléctrico integral de una instalación residencial
mediante simbología normalizada, criterios normativos, cuadro de cargas,
selección de conductores, protecciones y puesta a tierra.

\subsection{Producto mínimo}\label{producto-muxednimo}

\begin{itemize}
\tightlist
\item
  Plano arquitectónico con instalación eléctrica.
\item
  Cuadro de cargas.
\item
  Diagrama unifilar.
\item
  Cuadro de simbología.
\item
  Criterio de selección de conductores y protecciones.
\item
  Identificación de puesta a tierra.
\item
  Informe técnico.
\end{itemize}

\subsection{Criterios técnicos
básicos}\label{criterios-tuxe9cnicos-buxe1sicos}

\begin{itemize}
\tightlist
\item
  Separación de circuitos de iluminación, tomacorrientes y cargas
  especiales.
\item
  Uso de simbología normalizada.
\item
  Protección termomagnética por circuito.
\item
  Consideración de conductor de protección.
\item
  Ubicación funcional del tablero.
\end{itemize}

\section{Práctica 2: Diseño de plano eléctrico
industrial}\label{pruxe1ctica-2-diseuxf1o-de-plano-eluxe9ctrico-industrial}

\subsection{Objetivo}\label{objetivo-1}

Diseñar e interpretar una instalación eléctrica industrial básica,
diferenciando circuitos de potencia y control, diagramas unifilares,
multifilares, protecciones, cargas e instrumentos de medición.

\subsection{Producto mínimo}\label{producto-muxednimo-1}

\begin{itemize}
\tightlist
\item
  Plano o esquema de instalación industrial.
\item
  Diagrama de potencia.
\item
  Diagrama de control.
\item
  Identificación de instrumentos de medición.
\item
  Análisis de magnitudes eléctricas.
\item
  Informe técnico.
\end{itemize}

\section{Práctica 3: Fundamentos de programación con
PLCs}\label{pruxe1ctica-3-fundamentos-de-programaciuxf3n-con-plcs}

\subsection{Objetivo}\label{objetivo-2}

Configurar, conectar y programar manualmente un micro-PLC Siemens LOGO!
mediante funciones de temporización para el control secuencial de cargas
eléctricas.

\subsection{Producto mínimo}\label{producto-muxednimo-2}

\begin{itemize}
\tightlist
\item
  Diagrama de conexión.
\item
  Programa o lógica de control.
\item
  Evidencia de funcionamiento.
\item
  Tabla de entradas y salidas.
\item
  Informe técnico.
\end{itemize}

\subsection{Tabla de entradas y salidas
sugerida}\label{tabla-de-entradas-y-salidas-sugerida}

{\def\LTcaptype{none} % do not increment counter
\begin{longtable}[]{@{}lll@{}}
\toprule\noalign{}
Variable & Dirección & Descripción \\
\midrule\noalign{}
\endhead
\bottomrule\noalign{}
\endlastfoot
Marcha & I1 & Pulsador normalmente abierto \\
Paro / Reset & I2 & Pulsador normalmente cerrado \\
Carga & Q1 & Lámpara piloto o bobina de contactor \\
Temporizador & TON & Retardo a la conexión de 5 s \\
\end{longtable}
}

\bookmarksetup{startatroot}

\chapter{Evaluación}\label{evaluaciuxf3n}

\section{Distribución general}\label{distribuciuxf3n-general}

{\def\LTcaptype{none} % do not increment counter
\begin{longtable}[]{@{}lr@{}}
\toprule\noalign{}
Componente & Valoración \\
\midrule\noalign{}
\endhead
\bottomrule\noalign{}
\endlastfoot
Aprendizaje en contacto con el docente & 30 puntos \\
Aprendizaje autónomo & 20 puntos \\
Aprendizaje práctico-experimental & 20 puntos \\
Evaluación interciclo & 10 puntos \\
Evaluación final & 20 puntos \\
\textbf{Total} & \textbf{100 puntos} \\
\end{longtable}
}

\section{Actividades evaluativas
sugeridas}\label{actividades-evaluativas-sugeridas}

{\def\LTcaptype{none} % do not increment counter
\begin{longtable}[]{@{}llllr@{}}
\toprule\noalign{}
Actividad & Bloque & Técnica & Instrumento & Puntaje referencial \\
\midrule\noalign{}
\endhead
\bottomrule\noalign{}
\endlastfoot
Evaluación de conocimientos & 1 & Prueba & Cuestionario & 10 \\
Webinar & 2 & Informe / ensayo & Rúbrica & 5 \\
Investigación formativa & 2 & Ensayo & Rúbrica & 10 \\
Evaluación de conocimientos & 3 & Prueba & Cuestionario & 5 \\
Práctica residencial & 1 & Informe & Rúbrica & 20 \\
Práctica industrial & 2 & Informe & Rúbrica & 20 \\
Práctica PLC & 3 & Informe & Rúbrica & 20 \\
Proyecto final & 3 & Proyecto & Rúbrica & 10 \\
\end{longtable}
}

\section{Rúbrica general para
prácticas}\label{ruxfabrica-general-para-pruxe1cticas}

{\def\LTcaptype{none} % do not increment counter
\begin{longtable}[]{@{}
  >{\raggedright\arraybackslash}p{(\linewidth - 8\tabcolsep) * \real{0.1875}}
  >{\raggedright\arraybackslash}p{(\linewidth - 8\tabcolsep) * \real{0.1875}}
  >{\raggedright\arraybackslash}p{(\linewidth - 8\tabcolsep) * \real{0.1875}}
  >{\raggedright\arraybackslash}p{(\linewidth - 8\tabcolsep) * \real{0.1875}}
  >{\raggedleft\arraybackslash}p{(\linewidth - 8\tabcolsep) * \real{0.2500}}@{}}
\toprule\noalign{}
\begin{minipage}[b]{\linewidth}\raggedright
Criterio
\end{minipage} & \begin{minipage}[b]{\linewidth}\raggedright
Excelente
\end{minipage} & \begin{minipage}[b]{\linewidth}\raggedright
Bueno
\end{minipage} & \begin{minipage}[b]{\linewidth}\raggedright
Básico
\end{minipage} & \begin{minipage}[b]{\linewidth}\raggedleft
Puntaje
\end{minipage} \\
\midrule\noalign{}
\endhead
\bottomrule\noalign{}
\endlastfoot
Cumplimiento técnico & Aplica normativa, cálculos y criterios de
seguridad de forma consistente & Presenta criterios técnicos con errores
menores & Presenta criterios incompletos & 4 \\
Representación gráfica & Plano/diagrama claro, ordenado y normalizado &
Plano/diagrama comprensible con detalles por corregir & Plano/diagrama
incompleto & 4 \\
Cálculos y justificación & Cálculos completos y correctamente
interpretados & Cálculos parcialmente correctos & Cálculos insuficientes
& 4 \\
Informe técnico & Estructura clara, redacción técnica y evidencias &
Presenta estructura aceptable & Informe débil o desordenado & 4 \\
Seguridad y procedimiento & Identifica riesgos y aplica prevención &
Reconoce parcialmente riesgos & Omite criterios de seguridad & 4 \\
\textbf{Total} & & & & \textbf{20} \\
\end{longtable}
}

\bookmarksetup{startatroot}

\chapter{Proyecto final del módulo}\label{proyecto-final-del-muxf3dulo}

\section{Título sugerido}\label{tuxedtulo-sugerido}

\textbf{Diseño, interpretación y automatización básica de una
instalación eléctrica civil o industrial}

\section{Propósito}\label{propuxf3sito}

Integrar los contenidos de normalización, simbología, seguridad,
diagramas eléctricos, mediciones e introducción a la automatización
industrial mediante un caso aplicado.

\section{Entregables}\label{entregables}

\begin{enumerate}
\def\labelenumi{\arabic{enumi}.}
\tightlist
\item
  Descripción del caso de estudio.
\item
  Plano eléctrico o esquema de instalación.
\item
  Cuadro de cargas.
\item
  Diagrama unifilar.
\item
  Diagrama de control y potencia, si aplica.
\item
  Selección preliminar de conductores y protecciones.
\item
  Propuesta de mediciones eléctricas.
\item
  Identificación de riesgos eléctricos.
\item
  Propuesta básica de automatización con PLC o relés.
\item
  Informe técnico final.
\end{enumerate}

\section{Estructura del informe}\label{estructura-del-informe}

\begin{enumerate}
\def\labelenumi{\arabic{enumi}.}
\tightlist
\item
  Portada.
\item
  Introducción.
\item
  Objetivos.
\item
  Marco normativo y técnico.
\item
  Desarrollo del diseño.
\item
  Cálculos.
\item
  Diagramas.
\item
  Análisis de seguridad.
\item
  Automatización propuesta.
\item
  Conclusiones.
\item
  Bibliografía.
\item
  Anexos.
\end{enumerate}

\bookmarksetup{startatroot}

\chapter{Referencias}\label{referencias}

\protect\phantomsection\label{refs}




\end{document}
